\scanpage{9}
\mspage{8}
\mspage{IV-965}
\uline{Corollaire} (18.9.4) .- \uline{Soit} $A$ \uline{un anneau local noethérien, hensélien et dont les fibres formelles sont géométriquement normales} \deleted{\illegible{1}} (par exemple un anneau local noethérien hensélien et \uline{excellent}), \uline{et posons} $A'=\hat A$. \uline{Soient} $X$ \uline{un} $A$-\uline{préschéma}, $X'=X\otimes_A A'$, $g:X'\to X$ \uline{la projection canonique. Alors l'application} $U\rightsquigarrow g^{-1}(U)$ \uline{est une bijection de l'ensemble des parties à la fois ouvertes et fermées de} $X$ \uline{sur l'ensemble des parties à la fois ouvertes et fermées de} $X'$.

Comme $g$ est un morphisme \deleted{\illegible{1}} fidèlement plat et quasi-compact, on sait (2.3.12) que la topologie de $X$ est \uline{quotient} de celle de $X'$ par la relation d'équivalence définie par $f$. Tout revient donc à voir \deleted{\unsure{que}} qu'une partie ouverte et fermée $U'$ de $X'$ est \uline{saturée} pour cette relation. Or, comme le morphisme $\operatorname{Spec}(A')\to\operatorname{Spec}(A)$ \deleted{\illegible{1}} a ses fibres géométriquement connexes (18.9.1), les fibres de $g$ sont \uline{connexes}, d'où aussitôt notre assertion.

En particulier, si $X$ est localement connexe (ce qui sera le cas si $X$ est localement noethérien), alors, pour toute \deleted{\unsure{partie}} composante connexe $U$ de $X$ (qui est ouverte et fermée dans $X$), $g^{-1}(U)$ est connexe (Bourbaki, \uline{Top.gén.}, chap.I, $3^e$ éd., \S\ 11, $n^\circ3$, \inserted{prop 7}). Si l'on désigne par $\pi_0(X)$ l'ensemble des composantes connexes de $X$, , l'application $\pi_0(X')\to\pi_0(X)$ déduite canoniquement de $g$ est donc \uline{bijective}.

\margin{``séparés'' inutile !}
\deleted{
\uline{Corollaire} (18.9.5) .- \uline{Sous les hypothèses de} (18.9.4), \uline{le foncteur}
\[
(18.9.5.1)\qquad Z\rightsquigarrow Z\otimes_A A'
\]
\uline{de la catégorie des} \deleted{$X$-}\inserted{pré}\uline{schémas étales} \inserted{\deleted{(resp. étales et finis)} (et séparés)} \uline{sur} $X$ \uline{dans celle des} \deleted{\illegible{1}}\inserted{pré}\uline{schémas étales} \inserted{\deleted{(resp. étales et finis)} (et séparés)} \uline{sur} $X'$, \uline{est pleinement fidèle}.

En effet, si $Z_1,Z_2$ soit deux $X$-schémas étales sur $X$, $\operatorname{Hom}_X(Z_1,Z_2)$ s'identifie à $\Gamma(Z_1\times_X Z_2/Z_1)$, et $Z_1\times_X Z_2$ est étale et séparé sur $Z_1$, donc finalement (17.9.3), $\operatorname{Hom}_X(Z_1,Z_2)$ s'identifie fonctoriellement à l'ensemble des parties ouvertes et fermées $U$ de $Z_1\times_X Z_2$ telles que la restriction de la projection $p:Z_1\times_X Z_2\to Z_1$ à $U$ soit un morphisme surjectif et radiciel. \unsure{Le|La} corollaire résulte donc aussitôt de (18.9.4) et du fait que, si $Z_{\unsure{1|i}}'=Z_{\unsure{1|i}}\otimes_A A'$ \deleted{\illegible{1}} $(\unsure{1|i}=1,2)$, on a $Z_1'\times_{X'}Z_2'=(Z_1\times_X Z_2)\otimes_A A'$.

\uline{Remarques} (18.9.6) .- (i) \deleted{Dans tous les cas connus} \inserted{il est possible que si le corps résiduel de $A$ est de car. nulle,} le foncteur (18.9.5.1) \deleted{est}\inserted{induit}
}
\margin{\unsure{le resp.} n'ajoute rien ; $X$-schéma risque d'induire en erreur---}
